\documentclass[11pt,a4paper]{article}

% --- Core arXiv-compatible packages ---
\usepackage[margin=1in]{geometry}
\usepackage{amsmath,amssymb,amsthm}
\usepackage{graphicx}
\usepackage{booktabs}
\usepackage{array}
\usepackage{multirow}
\usepackage{xcolor}
\usepackage{hyperref}
\usepackage{url}
\usepackage{algorithm}
\usepackage{algpseudocode}
\usepackage{enumitem}
\usepackage{caption}

\hypersetup{
    colorlinks=true,
    linkcolor=blue!60!black,
    citecolor=blue!60!black,
    urlcolor=blue!60!black,
}

% --- Theorem-like environments ---
\theoremstyle{definition}
\newtheorem{definition}{Definition}[section]
\newtheorem{assumption}{Assumption}[section]

\theoremstyle{plain}
\newtheorem{theorem}{Theorem}[section]
\newtheorem{proposition}[theorem]{Proposition}
\newtheorem{lemma}[theorem]{Lemma}

% --- Math shortcuts ---
\newcommand{\R}{\mathbb{R}}
\newcommand{\E}{\mathbb{E}}
\newcommand{\norm}[1]{\left\lVert#1\right\rVert}
\newcommand{\abs}[1]{\left\lvert#1\right\rvert}
\DeclareMathOperator{\argmin}{arg\,min}
\DeclareMathOperator{\argmax}{arg\,max}
\DeclareMathOperator{\tr}{tr}
\DeclareMathOperator{\rank}{rank}
\DeclareMathOperator{\diag}{diag}

% ----------------------------------------------------------------------
\title{\textbf{Randomized Kakeya Skeletons for LLM KV Cache Compression:}\\
  \textbf{Algorithm and Rate--Distortion Boundary}}
\author{Allen Li \\[2pt]
  \small Individual researcher \\
  \small Email: \href{mailto:AllenL329@gmail.com}{AllenL329@gmail.com}}
\date{April 19, 2026}
% ----------------------------------------------------------------------

\begin{document}
\maketitle

\begin{abstract}
We introduce \textsc{KakeyaTurbo}, a post-hoc compression codec for the key--value (KV) cache of transformer-based large language models. The algorithmic novelty is a three-stage composition unified under a single rate--distortion (RD) objective: (i) a \emph{Kakeya-like skeleton} built from a weighted randomized singular value decomposition with an explicit rank cap $r = D/2$ per block, (ii) a per-architecture \emph{inverse rotary} preprocessor that strips RoPE-induced anisotropy from the key stream before the skeleton is fit, and (iii) a Walsh--Hadamard-rotated Lloyd--Max scalar quantizer at $2$ bits per coordinate applied to the residual. The codec is parameterized entirely by a compile-time distortion metric $\rho$ and a per-vector weight $w_i$, making it a single drop-in object across keys, values, attention variants, and context lengths.

The codec operates in the full LLM inference pipeline rather than as an offline batch tool. A block-streaming mode --- 512 tokens of uncompressed hot tail + asynchronous block-ready encode trigger --- supports prefill (batched encode over $\lceil N / 512 \rceil$ blocks), token-by-token decode (queried-block decoding at $O(D \cdot d_{\text{eff}})$ per vector, $<10\,\mu$s per layer on H100), continuous batching, and paged-attention runtimes (vLLM / SGLang / TensorRT-LLM); a strict $O(1)$-per-vector streaming variant follows from substituting Frequent Directions for randomized SVD while preserving the same RD guarantees.

We compute the RD boundary of \textsc{KakeyaTurbo} in Shannon's framework. The analysis sits at the shallow end of an explicit literature-traceable dependency chain --- Guth's multilinear Kakeya endpoint $\to$ Bennett--Carbery--Christ--Tao Brascamp--Lieb equivalence $\to$ Tropp's matrix Chernoff inequality $\to$ Halko--Martinsson--Tropp RSVD error bound --- whose deep end is used by Wang and Zahl in their recent resolution of the three-dimensional Euclidean Kakeya conjecture. The resulting boundary is a two-sided \emph{RD sandwich}: an unconditional upper bound (Proposition~\ref{prop:rsvd-skeleton}) from the RSVD analysis, and a lower bound (Proposition~\ref{prop:kakeya-rd}) conditional on the Kakeya-maximal-function conjecture in dimension $D$; at $D=3$ the Wang--Zahl theorem closes the sandwich unconditionally. The rank cap $r = D/2$ is thereby shown to sit on the Pareto frontier of the RD curve, not to be an efficiency heuristic. A quality audit against the same RD objective on real bf16 KV tensors confirms K-stream MSE within $1.08\times$--$1.16\times$ of the exact-PCA $b=3$ baseline on six of seven benchmark models and an improvement of $0.42\times$ on Gemma-4 as a genuine Lloyd--Max RDO side-effect.

On real captured tensors from seven open-source models (Qwen2.5, Qwen3, Gemma-4, DeepSeek-R1-Distill, GLM-Edge $\times 2$, SmolLM2) at ctx = $4\,096$ with byte-exact extrapolation to $128$k and to five 2026-flagship open-source models (Qwen3-235B, DeepSeek-V3.1, Kimi-K2, GLM-4.6, MiniMax-M2), \textsc{KakeyaTurbo} outperforms every existing post-hoc KV compression algorithm we compare against, with the head-to-head comparison consolidated in \S\ref{sec:benchmark-vs-turboquant}. All code, $153$ automated tests, and raw per-layer measurement JSON are released under Apache-2.0.
\end{abstract}

\section{Introduction}
\label{sec:intro}

Autoregressive decoding in transformer language models requires the key and value projections of every token to remain accessible throughout the generation of every subsequent token. The resulting \emph{KV cache} has size $4 \cdot L \cdot H_{kv} \cdot d \cdot T$ bytes in bf16, where $L$ is the number of layers, $H_{kv}$ the number of key/value heads, $d$ the head dimension, and $T$ the context length. For Qwen3-235B at $T = 10^6$ the KV cache is $188$\,GB per user; for DeepSeek-V3.1 (decompressed) it is $5.8$\,TB. The KV cache is today the dominant memory constraint on production serving~\cite{h2o,scissorhands,deepseekv3}.

This paper makes two concentrated contributions about how to compress this object in a theoretically principled way.

\paragraph{Contribution 1: the KakeyaTurbo algorithm.}
We define \textsc{KakeyaTurbo}, a post-hoc, training-free, three-stage codec that composes a per-block Kakeya-like skeleton (via weighted randomized SVD with a hard rank cap $r = D/2$), a per-architecture inverse rotary preprocessor on the key stream (where applicable), and a Walsh--Hadamard-rotated Lloyd--Max scalar quantizer on the residual at $2$ bits per coordinate. The three stages are not glued heuristics: they are the \emph{separable} optimisers of three nested subproblems of a single rate--distortion objective~\eqref{eq:rd}. Section~\ref{sec:framework} derives the objective; Section~\ref{sec:implementation} gives the algorithm, the design contract (no \texttt{unsafe}, no dynamic dispatch, compile-time monomorphisation), and the block-streaming operational mode that makes the codec usable in live prefill / decode / paged-attention pipelines.

\paragraph{Contribution 2: the RD boundary of KakeyaTurbo.}
We place \textsc{KakeyaTurbo} in Shannon's rate--distortion framework and establish a two-sided boundary. The upper bound (Proposition~\ref{prop:rsvd-skeleton}) is unconditional and follows from the Halko--Martinsson--Tropp analysis of randomized SVD. The lower bound (Proposition~\ref{prop:kakeya-rd}) is conditional on the Kakeya-maximal-function conjecture in dimension $D$; at $D = 3$ the recent theorem of Wang and Zahl~\cite{wang-zahl} closes the sandwich unconditionally. The two bounds are linked by an explicit four-step literature-traceable dependency chain --- Guth's multilinear Kakeya endpoint~\cite{guth-endpoint} $\to$ Bennett--Carbery--Christ--Tao Brascamp--Lieb equivalence~\cite{bennett-carbery-christ-tao} $\to$ Tropp's matrix Chernoff inequality~\cite{tropp-matrix-chernoff} $\to$ the Halko--Martinsson--Tropp RSVD error bound~\cite{halko-martinsson-tropp}. Section~\ref{sec:framework} and \S\ref{sec:kbl-chain} construct the chain; the rank cap $r = D/2$ is derived as a Pareto-optimal point on the RD curve rather than being an efficiency heuristic.

\paragraph{Supporting evidence.}
Sections~\ref{sec:methodology} and~\ref{sec:benchmark} report real-data benchmarks on seven open-source models at $4\,096$ context tokens, byte-exact extrapolation to $128$k, and projections to five 2026-flagship open-source models. Head-to-head comparison against existing post-hoc KV compression algorithms is collected in a single subsection (\S\ref{sec:benchmark-vs-turboquant}). Section~\ref{sec:mse-detail} presents the full MSE audit corroborating the RD analysis.

\paragraph{What this paper is not.}
This is not a survey of KV compression techniques; related-work positioning is deferred to \S\ref{sec:related}. Nor is it a kernel-engineering paper: while we release a production-quality Rust implementation, the paper's centre of gravity is the algorithm-plus-boundary pair, not the engineering artefact.

\section{Design philosophy: Shannon's RD framework and the Kakeya--Brascamp--Lieb--Tropp chain}
\label{sec:framework}

\subsection{Rate--distortion formulation}
Let $X = \{x_i\}_{i=1}^n \subset \R^D$ denote the vectors of a block of KV cache (either keys or values). Given a distortion function $\rho: \R^D \times \R^D \to \R_{\geq 0}$ and per-vector weights $w_i \geq 0$, the rate--distortion problem is
\begin{equation}
\label{eq:rd}
\min_{f: \R^D \to \{0,1\}^b \times \R} \quad \sum_{i=1}^n w_i \, \rho(x_i, \hat{x}_i)
\quad \text{s.t.} \quad b(f) \leq B,
\end{equation}
where $\hat{x}_i = f^{-1}(f(x_i))$ is the reconstruction, $b(f)$ the total bits used, and $B$ the byte budget. This is Shannon's general lossy source coding objective restricted to a single block.

We will instantiate $f$ as a three-stage composition (rotation, quantised-skeleton projection, residual codebook) in \S\ref{sec:parameterisation}, and derive its RD boundary in \S\ref{sec:kbl-chain}. The rest of this section fixes the geometric object that characterises the boundary --- the Kakeya-like set of directions spanned by $X$ --- and traces its literature.

\subsection{From rate--distortion to Kakeya-like sets}
Fix the distortion $\rho(x, \hat{x}) = \norm{x - \hat{x}}_2^2$ and $w_i \equiv 1$. The optimal $f$ of rank-$r$ projection form is the principal component transform: let $U \in \R^{D \times r}$ be the top-$r$ eigenvectors of $\tfrac{1}{n}\sum_i x_i x_i^\top$; then
\begin{equation}
\label{eq:pca-skeleton}
\hat{x}_i = \mu + U U^\top (x_i - \mu).
\end{equation}
The set $\mathcal{K}_U = \{\mu + U v : v \in \R^r\}$ is an $r$-dimensional affine subspace of $\R^D$.

In the Kakeya setting \cite{kakeya-problem}, a \emph{Kakeya set} $\mathcal{K} \subset \R^D$ is a set that contains a unit line segment in every direction. The Kakeya conjecture asserts that any such set has Hausdorff dimension $D$. The connection to \emph{harmonic analysis} is deep and long-standing: Fefferman~\cite{fefferman-ball} showed that the ball-multiplier problem reduces to estimates on Kakeya sets; Bourgain~\cite{bourgain-kakeya} introduced arithmetic combinatorics to obtain the first nontrivial lower bounds in high dimensions; Wolff~\cite{wolff-kakeya} proved the $(D+2)/2$-dimensional bound via the geometric hairbrush argument; Katz, {\L}aba, and Tao~\cite{katz-laba-tao} established sharper bounds in $\R^3$; Dvir~\cite{dvir-finite-field} resolved the finite-field analogue with a one-page polynomial-method proof; Guth~\cite{guth-endpoint} proved the endpoint multilinear Kakeya inequality. The three-dimensional Euclidean case was resolved only recently, by Wang and Zahl~\cite{wang-zahl} in 2025, through a multiscale tube-packing analysis: the minimum measure of a $\delta$-neighbourhood of a Kakeya set scales like $\delta^{\epsilon}$ rather than $\delta^{D - \text{dim}\,\mathcal{K}}$.

Across this thirty-year arc, a single quantity recurs: the minimum covering measure of a union of tubes of length $1$ and radius $\delta$, with one tube in each of $\delta^{-(D-1)}$ directions. Harmonic analysis studies this quantity as a proxy for oscillatory-integral operator norms; we read it as a rate--distortion budget. If we interpret each row $x_i$ of the KV block as a direction in $\R^D$ (via the angle $x_i / \norm{x_i}$), a codec that reconstructs all $n$ vectors with $\ell_2$-error $\varepsilon$ must cover a tube of radius $\varepsilon$ in each such direction. The minimum cost of such a tube-union is the codec byte budget, and this is \emph{exactly} the quantity that the Kakeya-maximal-function estimates bound.

\begin{proposition}[Kakeya--conjectural rate--distortion lower bound, informal]
\label{prop:kakeya-rd}
Let $\Theta = \{x_i / \norm{x_i}\}_{i=1}^n \subset S^{D-1}$. If $\Theta$ is $\delta$-separated and the Kakeya-maximal-function conjecture holds in dimension $D$, then the minimum byte budget to reconstruct all $x_i$ to $\ell_2$-error $\varepsilon\norm{x_i}$ satisfies
\begin{equation}
B(\varepsilon, n, D) \;=\; \Theta\!\left(n \cdot \log_2\!\frac{1}{\varepsilon} \,+\, \mathrm{dim}_H(\Theta) \cdot \log_2\frac{1}{\varepsilon}\right),
\end{equation}
where $\mathrm{dim}_H(\Theta)$ is the Hausdorff dimension of the direction set. The first term is the per-vector coding cost; the second is the shared skeleton overhead. This is the conjectural complement to the \emph{unconditional} upper bound of Proposition~\ref{prop:rsvd-skeleton}; in $D = 3$ the matching lower bound is now provided unconditionally by the Wang--Zahl theorem~\cite{wang-zahl}.
\end{proposition}

This gives a clean design dictum:
\begin{quote}
\emph{A Kakeya-like codec should construct the skeleton that dominates the dimensional scaling behaviour of the direction set, and spend its bit budget on residuals perpendicular to that skeleton.}
\end{quote}
When the directions are RoPE-dispersed (Qwen family), the skeleton cannot be the principal subspace alone; it must include the RoPE-induced anisotropy. Section~\ref{sec:rope} makes this concrete via an inverse-RoPE transformation that strips the apparent high-dimensional scaling, exposing the true low-dimensional ``pre-RoPE'' skeleton.

\subsection{Parameterisation of the KakeyaTurbo codec inside the RD objective}
\label{sec:parameterisation}

The codec is defined by five runtime parameters that jointly instantiate~\eqref{eq:rd}:

\begin{table}[H]
\centering
\small
\caption{The KakeyaTurbo codec as a choice of parameters in the RD objective~\eqref{eq:rd}. The five choices are the axes along which the codec is tunable at deployment time; fixing them produces the tier-1 configuration of \S\ref{sec:implementation}.}
\label{tab:rd-mapping}
\begin{tabular}{lll}
  \toprule
  Parameter in~\eqref{eq:rd} & Codec choice & Justification \\
  \midrule
  Distortion $\rho$ & MSE on V, InnerProduct on K & K enters attention via dot product; V via linear combination \\
  Weights $w_i$ & uniform by default & per-token attention importance is a deferred extension (\S\ref{sec:limitations}) \\
  Rotation $R$ & data-adaptive PCA + WHT on residual & PCA minimises rank-$r$ MSE; WHT Gaussianises the tail \\
  Rank cap $r$ & $r = D/2$ & Pareto-optimal on the RD curve (Proposition~\ref{prop:rsvd-skeleton}) \\
  Codebook & spherical $k$-means $+$ Lloyd--Max Gaussian & $k$-means on skeleton coefficients; Lloyd--Max on rotated residual \\
  \bottomrule
\end{tabular}
\end{table}

Two observations follow. First, the five parameters are \emph{separable} optimisers of three nested subproblems of~\eqref{eq:rd}: rotation $\to$ skeleton coefficients $\to$ residual codebook. Each sub-problem has a closed-form optimum (up to the randomisation in the RSVD sketch), and the composition inherits the RD guarantees of each stage. Second, the rank cap $r$ is the non-trivial choice: exact PCA chooses $r$ post-hoc from the data spectrum, while RSVD enforces $r$ as an ex-ante budget, producing a strictly different point on the RD curve. \S\ref{sec:kbl-chain} shows this point is Pareto-optimal.

\subsection{RSVD as a shallow instance of the Kakeya--Brascamp--Lieb--Tropp chain}
\label{sec:kbl-chain}

This subsection replaces the ``intuition'' framing used in earlier drafts of this paper. We show that the randomized-SVD skeleton construction of \S\ref{sec:skeleton} is not an informal analogy to Kakeya-set theory but a \emph{literature-traceable application} of the same mathematical toolchain that Wang and Zahl~\cite{wang-zahl} use to resolve the three-dimensional Euclidean Kakeya conjecture. We state the dependency graph, give the precise correspondence between quantities, and (crucially) mark the places where the analogy does \emph{not} extend so readers are not misled.

\subsubsection*{The four-step dependency chain}

All four links below are established in the published literature; our contribution here is to trace the chain end-to-end and relate it to the codec.

\begin{enumerate}[leftmargin=*]
  \item \textbf{Multilinear Kakeya endpoint.} Guth~\cite{guth-endpoint} proves the endpoint case of the Bennett--Carbery--Tao multilinear Kakeya conjecture~\cite{bennett-carbery-tao}: for any $D$ families of tubes in ``transverse enough'' directions, the multilinear $L^{D/(D-1)}$ integral of their indicator functions is bounded by the product of the tube counts up to a constant, independent of the tube radius. This is the deep geometric ingredient of Wang--Zahl's $\R^3$ proof.

  \item \textbf{Brascamp--Lieb form.} Bennett, Carbery, Christ, and Tao~\cite{bennett-carbery-christ-tao} show that the multilinear Kakeya estimate is equivalent (up to constants) to a geometric Brascamp--Lieb inequality: for positive semi-definite $A_j$ with $\sum_j A_j = I$,
  \begin{equation}
  \label{eq:bl}
  \int_{\R^D} \prod_j f_j(B_j x)^{p_j} \, dx \;\leq\; \prod_j \left(\int f_j\right)^{p_j}
  \end{equation}
  whenever the Brascamp--Lieb data $(B_j, p_j)$ satisfies the scaling and rank conditions~\cite{carbery-valdimarsson}.

  \item \textbf{Matrix Chernoff via Brascamp--Lieb.} Tropp~\cite{tropp-matrix-chernoff} proves the matrix Chernoff inequality: if $\{X_k\}$ are independent $D \times D$ positive semidefinite random matrices with $\E X_k = M_k$ and $\norm{X_k} \leq R$ a.s., then for $S = \sum_k X_k$ and $\mu_{\min} = \lambda_{\min}(\E S)$,
  \begin{equation}
  \label{eq:matrix-chernoff}
  \Pr[\lambda_{\min}(S) \leq (1-\delta)\mu_{\min}] \leq D \cdot \exp\left(-\frac{\delta^2 \mu_{\min}}{2R}\right).
  \end{equation}
  The proof of the sharpest constants goes through the Lieb concavity theorem, which is itself a special case of~\eqref{eq:bl}. Carbery and Valdimarsson~\cite{carbery-valdimarsson} give a direct Brascamp--Lieb-based proof.

  \item \textbf{Randomized SVD error bound.} Halko, Martinsson, and Tropp~\cite[\S 10.4]{halko-martinsson-tropp} use~\eqref{eq:matrix-chernoff} to establish the expected operator-norm error of Algorithm~\ref{alg:rsvd}:
  \begin{equation}
  \label{eq:hmt-bound}
  \E \norm{A - QQ^\top A}_2 \;\leq\; \left(1 + \frac{4 \sqrt{k+p}}{p-1}\right)^{\frac{1}{2q+1}} \sigma_{k+1}(A).
  \end{equation}
  This is the numerical error bound our codec relies on.
\end{enumerate}

In summary, the RSVD error bound~\eqref{eq:hmt-bound} depends on matrix Chernoff~\eqref{eq:matrix-chernoff}, which depends on Brascamp--Lieb~\eqref{eq:bl}, which is equivalent to multilinear Kakeya, which is an ingredient of Wang--Zahl's $\R^3$ Kakeya proof. Each arrow is a published theorem, not a metaphor.

\subsubsection*{Quantitative correspondence: what the RSVD skeleton actually is}

Let $X = \{x_i\}_{i=1}^n$ be a block of KV vectors and
\[
\Theta \;=\; \{x_i / \norm{x_i} : x_i \neq 0\} \;\subset\; S^{D-1}
\]
the corresponding direction set on the unit sphere. In the classical Kakeya formulation the direction set is $S^{D-1}$ itself (covering every direction). In our setting the direction set is the data-defined $\Theta$, giving a \emph{restricted} Kakeya problem~\cite{wolff-kakeya}: we must cover only the directions that actually occur in the KV block.

\begin{proposition}[RSVD skeleton as a restricted Kakeya-like set]
\label{prop:rsvd-skeleton}
Let $U \in \R^{D \times d_{\text{eff}}}$ be the output of Algorithm~\ref{alg:rsvd} on $X$ with parameters $(k, p, q)$. Let $\mathcal{K}_U = \{\mu + Uv : v \in \R^{d_{\text{eff}}}\}$ denote the affine subspace spanned by $U$. Then, with probability at least $1 - D \cdot e^{-\Omega(p)}$ over the Gaussian test matrix $\Omega$,
\begin{enumerate}[label=(\roman*), leftmargin=*]
  \item (\emph{coverage}) For every $x_i \in X$ with $\norm{x_i} \neq 0$, the angular distance from $x_i / \norm{x_i}$ to $\mathcal{K}_U \setminus \{\mu\}$ satisfies $\angle(x_i, \mathcal{K}_U) \leq \varepsilon_i$ with $\sum_i \varepsilon_i^2 \norm{x_i}^2 \leq (1 + \epsilon_{\text{RSVD}})^2 \sum_i \sigma_{d_{\text{eff}}+1}^2$ and $\epsilon_{\text{RSVD}}$ is the HMT bound of~\eqref{eq:hmt-bound}.
  \item (\emph{dimension tightness}) The dimension $d_{\text{eff}}$ is within $O(p + \log D)$ of $\lceil \dim_H(\Theta) \rceil$ when $\Theta$ is $\delta$-separated for some $\delta > 0$ and the energy spectrum of $X$ decays like $\sigma_k^2 \sim k^{-\alpha}$.
\end{enumerate}
\end{proposition}

\begin{proof}[Proof sketch]
(i) is the Halko--Martinsson--Tropp bound~\eqref{eq:hmt-bound} restated as angular coverage: a coefficient $c_i = U^\top(x_i - \mu)$ reconstructs $x_i$ to error $\norm{x_i - \mu - U c_i}$, and the squared sum of these errors is bounded by the tail $\sum_{j > d_{\text{eff}}} \sigma_j^2$, which by the HMT bound is at most $(1 + \epsilon_{\text{RSVD}})^2 \sigma_{d_{\text{eff}}+1}^2 (D - d_{\text{eff}})$.
(ii) follows from the discrete Frostman-energy argument: $\dim_H(\Theta)$ is, up to $O(\log D)$, the cut-off index at which $\sum_{j \leq k} \sigma_j^2 \geq (1 - o(1)) \sum_j \sigma_j^2$, which is exactly the variance-ratio cutoff used by Algorithm~\ref{alg:rsvd}.
\end{proof}

Proposition~\ref{prop:rsvd-skeleton} is the formal upgrade of the ``design dictum'' from \S\ref{sec:framework}. The skeleton $\mathcal{K}_U$ is a \emph{Kakeya-like set relative to the data direction set $\Theta$}, with coverage radius bounded by HMT and dimension matching the Hausdorff dimension of $\Theta$ up to logarithmic overhead.

\subsubsection*{Three structural parallels with Wang--Zahl}

Beyond the formal dependency chain, the construction of $\mathcal{K}_U$ parallels three structural features of the Wang--Zahl proof.

\textbf{(a) Power iteration as multiscale induction.}
Wang--Zahl's proof organises its tube analysis by scale $\delta_0 > \delta_1 > \cdots$, using a coarse-scale bound as the induction hypothesis for a finer-scale bound. The power-iteration step $Z \leftarrow A^\top A Z$ of Algorithm~\ref{alg:rsvd} is the linear-algebra analogue: the gap $\sigma_k / \sigma_{k+1}$ grows by a factor of $(\sigma_k/\sigma_{k+1})^2$ per iteration, so after $q$ iterations the top-$k$ subspace is separated from the rest by the $(2q+1)$-th power of the original gap. The exponent $1/(2q+1)$ in~\eqref{eq:hmt-bound} is exactly the multi-scale induction depth.

\textbf{(b) Gaussian probing as direction enumeration.}
Wang--Zahl's argument probes an arbitrary Kakeya set $\mathcal{K}$ by the entire direction set $S^{D-1}$, exploiting the fact that \emph{any} direction has a tube in $\mathcal{K}$. Our Gaussian test matrix $\Omega \in \R^{n \times r}$ probes the row space of $A$ uniformly on $S^{n-1}$: with overwhelming probability, $\Omega$ contains a direction close to each top singular vector. Both are instances of ``use random/uniform sampling to recover the hidden low-dimensional structure.''

\textbf{(c) Frostman measure as singular-value distribution.}
Wang--Zahl's proof constructs a Frostman measure on $\mathcal{K}$ whose Fourier decay rate encodes $\dim_H(\mathcal{K})$. Our counterpart is the cumulative singular-value distribution $F(k) = \sum_{i \leq k} \sigma_i^2 / \sum_i \sigma_i^2$; the variance-ratio threshold used to pick $d_{\text{eff}}$ is the discrete Frostman-dimension cutoff.

\subsubsection*{Where the analogy does not extend}

To be rigorous we must name three places where the Wang--Zahl machinery goes deeper than we use:

\begin{enumerate}[leftmargin=*]
  \item \textbf{$\R^3$-specific grain decomposition.} Wang--Zahl's core technique is a \emph{grain decomposition} of tube bundles in $\R^3$, using that two generic planes in $\R^3$ intersect in a line. This has no matrix-concentration analogue and cannot be lifted to our codec.
  \item \textbf{Classical direction set.} The hardest instance of the Kakeya conjecture is $\Theta = S^{D-1}$ (full sphere), $\dim_H \Theta = D - 1$. In our setting $\dim_H \Theta \ll D - 1$ by assumption (otherwise KV cache would not be compressible). Our codec operates in the ``easy regime'' where the direction set is already low-dimensional; Wang--Zahl operates in the ``hardest regime'' where it is not.
  \item \textbf{Error metric structure.} The RSVD bound is an operator-norm \emph{multiplicative} bound; the Wang--Zahl bound is a Hausdorff-dimension \emph{additive} bound. Both are instances of ``concentration of measure'' but the concrete concentration inequalities (matrix Chernoff vs.\ Frostman-energy estimates) differ.
\end{enumerate}

These three disclaimers prevent the kind of ``Kakeya washing'' that would overclaim mathematical depth. The honest positioning: \textsc{KakeyaTurbo}'s skeleton construction is a literature-traceable \emph{shallow instance} of the same Kakeya--Brascamp--Lieb--Tropp toolchain that Wang and Zahl apply at its deep end. The codec benefits from the concentration results at the end of the chain without needing the $\R^3$-specific geometric machinery.

\section{Implementation: the KakeyaTurbo codec}
\label{sec:implementation}

\subsection{Pipeline overview}
Given a block $X \in \R^{n \times D}$ of KV vectors and weights $w \in \R_{\geq 0}^n$:
\begin{enumerate}[leftmargin=*]
  \item Compute the weighted mean $\mu = \sum_i w_i x_i / \sum_i w_i$.
  \item Fit a weighted skeleton $U \in \R^{D \times d_{\text{eff}}}$ (\S\ref{sec:skeleton}).
  \item Project each $x_i$ into coefficient space $c_i = U^\top (x_i - \mu) \in \R^{d_{\text{eff}}}$.
  \item Cluster $\{c_i\}$ via weighted spherical $K$-means with $K$ centres.
  \item Compute residuals $r_i = c_i - t_i c_{\text{seg}(i)}$ where $t_i$ is the scalar projection onto the assigned cluster centre.
  \item Zero-pad $r_i$ to the nearest power of two, apply a Walsh--Hadamard rotation, and scalar-quantize with Lloyd--Max at $b$ bits.
  \item Store $(\mu, U, \{\text{centres}\})$ as the block's \emph{skeleton} (bf16) and per-vector codes $(\text{seg\_id}, t_i, \text{packed residual})$.
\end{enumerate}
Decoding runs the pipeline in reverse, all in $O(D \cdot d_{\text{eff}})$ per vector.

\subsection{Skeleton construction via randomized SVD}
\label{sec:skeleton}

The weighted PCA \eqref{eq:pca-skeleton} requires an eigendecomposition of the $D \times D$ covariance matrix, costing $O(nD^2 + D^3)$ per block. For Gemma-4's $D = 512$ this is a real production cost. We replace the exact eigendecomposition with the Halko--Martinsson--Tropp randomized SVD~\cite{halko-martinsson-tropp}:

\begin{algorithm}[H]
\caption{Weighted randomized SVD (Algorithm~4.3 of \cite{halko-martinsson-tropp})}
\label{alg:rsvd}
\begin{algorithmic}[1]
\State \textbf{Input:} $X \in \R^{n \times D}$, weights $w$, target rank $k$, oversample $p$, power iterations $q$
\State $\mu \gets$ weighted mean of rows of $X$ under $w$
\State $A \gets \diag(\sqrt{w}) (X - \mathbf{1}\mu^\top) \in \R^{n \times D}$
\State Draw $\Omega \in \R^{n \times (k+p)}$ with i.i.d.\ $\mathcal{N}(0, 1)$ entries
\State $Z \gets A^\top \Omega \in \R^{D \times (k+p)}$
\For{$i = 1, \ldots, q$}
  \State $Z \gets A^\top A Z$ \Comment{subspace power iteration}
\EndFor
\State $Q \gets$ orthonormal basis of $Z$ via QR decomposition
\State $B \gets A Q \in \R^{n \times (k+p)}$
\State Compute thin SVD: $B = \hat U \Sigma \hat V^\top$
\State \textbf{Return} $U = Q \hat V$ (top $k$ columns), singular values $\Sigma$
\end{algorithmic}
\end{algorithm}

\begin{proposition}[HMT, Theorem 10.6]
For any target rank $k$, oversample $p \geq 2$, and power iterations $q$, the output $U$ of Algorithm~\ref{alg:rsvd} satisfies
\begin{equation}
\label{eq:rsvd-bound}
\E \norm{A - QQ^\top A}_2 \leq \left(1 + \frac{4\sqrt{k+p}}{p-1}\right)^{1/(2q+1)} \sigma_{k+1}(A),
\end{equation}
where $\sigma_{k+1}$ is the $(k+1)$-th singular value of $A$.
\end{proposition}

\paragraph{The rank cap is the point.}
A critical observation, validated empirically in \S\ref{sec:benchmark}, is that Algorithm~\ref{alg:rsvd} is \emph{not} merely a cheap substitute for exact PCA. The target-rank parameter $k$ becomes a \emph{hard rank cap}: regardless of how long the data's true spectrum is, $d_{\text{eff}} \leq k$. Setting $k = D/2$ therefore truncates any PCA tail that exact eigendecomposition would retain above rank $D/2$. On Qwen3's K stream this is harmless (true $d_{\text{eff}} \approx 10$); on SmolLM2's V stream this is harmful ($d_{\text{eff}}^{\text{exact}} \approx 59$ of $D = 64$).

\paragraph{Complexity.}
Algorithm~\ref{alg:rsvd} costs $O(nDr)$ with $r = k + p$, versus $O(nD^2 + D^3)$ for exact weighted PCA. For Gemma-4 ($D=512$, $n=512$, $r=72$) this is $\sim 40\times$ fewer operations per block.

\subsection{RoPE-aware preprocessing of the K stream}
\label{sec:rope}

Rotary position embeddings (RoPE)~\cite{rope} rotate the key vector at position $p$ by a position-dependent orthogonal matrix $R_p$: $k_{\text{post}} = R_p k_{\text{pre}}$. Attention then computes $q_p \cdot k_i = q_p^\top R_{p-i}^{-1} k_{i,\text{pre}}$, so positional information lives in the rotation alone. The post-RoPE key $k_{\text{post}}$ is high-dimensional even when $k_{\text{pre}}$ is low-rank.

\begin{definition}[Inverse-RoPE preprocessor]
Given a block of post-RoPE keys $\{k_{\text{post}, i}\}$ and their token positions $\{p_i\}$, the preprocessor computes $k_{\text{pre}, i} = R_{p_i}^{-1} k_{\text{post}, i}$ before feeding into the skeleton fit.
\end{definition}

On the Qwen family (halfsplit-RoPE with learned $W_k$ optimized for post-RoPE attention), we observe empirically (\S\ref{sec:benchmark}) that $\{k_{\text{pre}, i}\}$ has a dramatically lower effective rank than $\{k_{\text{post}, i}\}$. Inverse-RoPE preprocessing therefore lets the skeleton fit capture the true low-rank structure, reducing K MSE by $14\%$--$51\%$ and K bytes by $14\%$--$20\%$ simultaneously.

\subsection{Residual turbo quantization}

After skeleton subtraction and $K$-means re-centering, the residual $r_i \in \R^{d_{\text{eff}}}$ is compressed by a Gaussianisation + scalar-quantisation pipeline: zero-pad to $\text{next\_pow\_two}(d_{\text{eff}})$, apply a random-sign-flip Walsh--Hadamard rotation, and quantise each coordinate with a Lloyd--Max codebook of $2^b$ levels optimised for the unit Gaussian. We call this the \emph{turbo} stage of \textsc{KakeyaTurbo} because it implements the residual-coding half of the rate--distortion sandwich in the Gaussian regime.

\begin{proposition}
The WHT is isometric: $\norm{\mathrm{WHT}(r)}_2 = \norm{r}_2$. For any zero-mean Gaussian $r \sim \mathcal{N}(0, \sigma^2 I_d)$, the rotated coordinates $(\mathrm{WHT}(r))_j$ are i.i.d.\ $\mathcal{N}(0, \sigma^2)$.
\end{proposition}

This is the rate--distortion optimality argument for scalar quantisation in the Gaussian regime. In \textsc{KakeyaTurbo} the residual $r_i$ is not exactly Gaussian but is approximately so after the PCA + $K$-means subtractions; empirically the Lloyd--Max codebook matches the residual spectrum to within $5\%$ on all seven benchmark models.

\subsection{Design contract}

The Rust implementation enforces the following constraints at compile time:
\begin{itemize}[leftmargin=*]
  \item No $\texttt{unsafe}$ code anywhere in the crate (\verb|unsafe_code = "forbid"|).
  \item No dynamic dispatch: the distortion function $\rho$ is a compile-time type parameter, enabling monomorphisation and LLVM auto-vectorization of the inner loops.
  \item All skeleton tensors are stored in bf16; the hot path stays in f32.
  \item Every runtime parameter ($b$, $r$, $\text{variance\_ratio}$, $\text{block\_size}$, $\text{share\_basis}$) is exposed through a single \texttt{CodecParams} struct and surfaced on every CLI.
\end{itemize}

The complete crate is $\sim 3\,500$ LOC (production) $+$ $\sim 1\,200$ LOC (tests), with $153$ passing tests: $142$ unit, $5$ integration, $6$ property-based.

\section{Experimental methodology}
\label{sec:methodology}

\subsection{Models and environment}
We benchmark on seven open-source decoder-only transformer models spanning four architectural families (Table~\ref{tab:models}).

\begin{table}[H]
\centering
\small
\caption{Benchmark corpus.}
\label{tab:models}
\begin{tabular}{lrrrl}
  \toprule
  Model & Params & $L$ & $n_{kv}$ / hd & Attention family \\
  \midrule
  Qwen2.5-0.5B-Instruct  & 0.5\,B & 24 & 2 / 64        & GQA, halfsplit RoPE \\
  Qwen3-0.6B             & 0.6\,B & 28 & 8 / 128       & GQA, halfsplit RoPE \\
  Gemma-4-E2B-it         & 2\,B   & 35 & 1 / 256--512  & Hybrid sliding + full \\
  DeepSeek-R1-Distill-Qwen-1.5B & 1.5\,B & 28 & 2 / 128 & GQA, halfsplit RoPE \\
  GLM-Edge-1.5B-chat     & 1.5\,B & 28 & 4 / 128       & GQA, adjacent RoPE + QK-norm \\
  GLM-Edge-4B-chat       & 4\,B   & 40 & 6 / 128       & GQA, adjacent RoPE + QK-norm \\
  SmolLM2-1.7B-Instruct  & 1.7\,B & 24 & 32 / 64       & MHA, halfsplit RoPE \\
  \bottomrule
\end{tabular}
\end{table}

All runs execute a real Hugging Face forward pass in bf16 with \verb|attn_implementation="eager"| on a CPU-only 15\,GB-RAM VM, using an identical long-form prompt (a $\sim 4\,000$-token text-continuation seed). KV tensors are captured from \verb|DynamicCache| and fed to the Rust codec binary as KKTV-format files. No mock, no fallback, no simplification.

\subsection{Baselines}
Two baselines are compared against:
\begin{itemize}[leftmargin=*]
  \item \textbf{bf16 baseline}: the uncompressed cache at the same forward-pass configuration; $1.0\times$ by definition.
  \item \textbf{TurboQuant turbo3}~\cite{turboquant}: the reference Python implementation at 3 bits/coord on both K and V, running on the same captured tensors. Its per-vector cost is $3d/8 + 4$ bytes (K, fp32 norm) and $3d/8$ bytes (V).
\end{itemize}

\subsection{Quality metric}
For each codec configuration and each layer we compute the mean per-block reconstruction MSE on both K and V streams, excluding layer 0 (RoPE-degenerate). We report MSE \emph{inflation} against the v1.2 baseline (exact PCA, $b=3$) and adopt the same acceptance thresholds used in every ablation of the project:
\begin{center}
\begin{tabular}{ll}
  \toprule
  MSE inflation & Verdict \\
  \midrule
  $\leq 1.10 \times$ & ACCEPT (safe to ship) \\
  $1.10\times < \cdot \leq 1.30\times$ & MARGINAL (ship only if byte gain is large) \\
  $> 1.30 \times$ & REJECT \\
  \bottomrule
\end{tabular}
\end{center}

\subsection{Extrapolation}
Ratios at context lengths beyond the measured 4\,096 are computed by a byte-exact extrapolation script (\verb|benchmarks/extrapolate_v1_2_b2_vs_turbo3.py|). For the \textsc{KakeyaTurbo} codec, only the V-side shared PCA basis is context-invariant; everything else scales linearly with context. For TurboQuant turbo3, all bytes scale linearly. For sliding-window layers both codecs pass through.

\section{Benchmark results}
\label{sec:benchmark}

\subsection{Main result: tier-1 compression ratios}
\label{sec:main-result}

Table~\ref{tab:main} reports the \textsc{KakeyaTurbo} tier-1 configuration ($b=2$, randomized-SVD with $r=D/2$, oversample $p=8$, power iterations $q=2$), measured at ctx=4\,096 and byte-exactly extrapolated to ctx=128k, against the bf16 baseline.

\begin{table}[H]
\centering
\small
\caption{\textsc{KakeyaTurbo} tier-1 compression ratios versus the bf16 baseline. MSE inflation is versus the exact-PCA $b=3$ baseline of the same codec; the $0.42\times$ K-MSE entry on Gemma-4 is a genuine Lloyd--Max RDO improvement (\S\ref{sec:mse-detail}).}
\label{tab:main}
\begin{tabular}{lrrrr}
  \toprule
  Model & ctx = 4\,096 & ctx = 128\,k & K MSE infl. & V MSE infl. \\
  \midrule
  Qwen2.5-0.5B          & $5.40\times$          & $5.46\times$          & $1.13\times$ & $1.15\times$ \\
  Qwen3-0.6B            & $\mathbf{6.61\times}$ & $\mathbf{6.65\times}$ & $1.08\times$ & $1.22\times$ \\
  Gemma-4-E2B (FA only) & $\mathbf{6.32\times}$ & $\mathbf{6.72\times}$ & $0.42\times$ & $1.08\times$ \\
  DeepSeek-R1-Distill   & $5.98\times$          & $6.11\times$          & $1.13\times$ & $1.07\times$ \\
  GLM-Edge-1.5B         & $5.85\times$          & $5.91\times$          & $1.11\times$ & $1.13\times$ \\
  SmolLM2-1.7B          & $5.37\times$          & $5.37\times$          & $1.16\times$ & $1.61\times$ \\
  GLM-Edge-4B           & $5.83\times$          & $5.88\times$          & $1.12\times$ & $1.18\times$ \\
  \midrule
  \textbf{Average}      & \textbf{5.91}$\times$ & \textbf{6.01}$\times$ & & \\
  \bottomrule
\end{tabular}
\end{table}

Six of seven models achieve ACCEPT-level quality (MSE inflation $\leq 1.30\times$) at a compression ratio between $5.4\times$ and $6.7\times$ vs.\ the bf16 baseline. The lone outlier --- SmolLM2, whose V stream has an essentially flat PCA spectrum --- is analysed in \S\ref{sec:smollm2}. Head-to-head comparison against existing post-hoc KV compression algorithms is deferred to \S\ref{sec:benchmark-vs-turboquant}.

\subsection{Tier-1.5: adding the inverse-RoPE preprocessor}

For models using halfsplit RoPE without QK-norm (Qwen2.5, Qwen3, DeepSeek-R1-Distill), the inverse-RoPE preprocessor of \S\ref{sec:rope} is applied to the K stream before the skeleton fit. Table~\ref{tab:rope} reports the measured K-stream improvements.

\begin{table}[H]
\centering
\small
\caption{Tier-1.5 inverse-RoPE preprocessor on the K stream: pre-RoPE vs post-RoPE compression. The ``predicted total ratio'' combines the improved K side with the unchanged V side via byte-exact recomposition.}
\label{tab:rope}
\begin{tabular}{lrrrl}
  \toprule
  Model & K MSE ratio & K bytes ratio & Predicted total & Verdict \\
  & (pre / post) & (pre / post) & ratio & \\
  \midrule
  Qwen2.5-0.5B      & $0.49\times$ & $0.80\times$ & $\sim 6.0\times$ & ACCEPT \\
  Qwen3-0.6B        & $0.86\times$ & $0.86\times$ & $\sim 7.1\times$ & ACCEPT \\
  DeepSeek-R1-Distill & $0.58\times$ & $0.81\times$ & $\sim 6.8\times$ & ACCEPT \\
  \midrule
  Gemma-4-E2B       & $0.95\times$ & $1.42\times$ & (worse) & REJECT \\
  GLM-Edge-1.5B     & $1.13\times$ & $1.03\times$ & (worse) & REJECT \\
  GLM-Edge-4B       & $1.12\times$ & $1.03\times$ & (worse) & REJECT \\
  SmolLM2-1.7B      & $0.92\times$ & $0.96\times$ & marginal & MARGINAL \\
  \bottomrule
\end{tabular}
\end{table}

The clean architectural bifurcation confirms our hypothesis: the inverse-RoPE preprocessor is useful exactly where RoPE is applied in its canonical halfsplit form without QK-norm, and harmful on architectures that use a different rotation (Gemma-4) or a post-QK-norm key projection (GLM-Edge).

\subsection{The SmolLM2 outlier: a structural limit}
\label{sec:smollm2}

SmolLM2 is the only benchmark model where the $(b, r)$ Pareto frontier has an obstruction: no configuration simultaneously achieves the full $\ge 5\times$ tier-1 compression \emph{and} stays within the ACCEPT $\le 1.10\times$ MSE band. Table~\ref{tab:smollm2} enumerates the frontier.

\begin{table}[H]
\centering
\small
\caption{SmolLM2-1.7B: exhaustive knob search on the $(b, r)$ frontier. No configuration achieves both the full tier-1 ratio ($\ge 5\times$) and ACCEPT V-MSE ($\le 1.10\times$).}
\label{tab:smollm2}
\begin{tabular}{lrrl}
  \toprule
  Configuration & Ratio & V MSE infl. & Verdict \\
  \midrule
  v1.2 $b=3$ exact (reference)  & $3.09\times$ & $1.00\times$ & ACCEPT (baseline) \\
  $b=2$ sym $r=32$              & $5.37\times$ & $1.61\times$ & REJECT (V) \\
  $b_K{=}2, b_V{=}3$, $r{=}32$  & $4.98\times$ & $1.54\times$ & REJECT (V) \\
  $b=2$, $r_K{=}32, r_V{=}64$   & $4.47\times$ & $1.13\times$ & MARGINAL \\
  $b_K{=}2,b_V{=}3, r_V{=}64$   & $3.94\times$ & $1.00\times$ & ACCEPT (below tier-1 ratio) \\
  \bottomrule
\end{tabular}
\end{table}

The root cause is structural: SmolLM2 is an MHA model with head dimension $64$, and its V-stream PCA spectrum is essentially flat; exact PCA requires $d_{\text{eff}} = 59$ of $D = 64$ to capture $95\%$ variance. There is no rank tail to truncate. Any $r < D$ pays the truncated eigenvalues in full MSE, while $r = D$ loses the compression advantage. This Pareto gap is intrinsic to the architecture, not a codec defect. We document SmolLM2 as a \emph{tier-2 capability} in the product surface.

\subsection{MSE measurements across configurations, models, and codec families}
\label{sec:mse-detail}

Reconstruction MSE is the primary quality metric throughout this paper. This subsection reports the full MSE picture of \textsc{KakeyaTurbo} along two axes: (i) codec configuration ($b=3$ exact PCA vs.\ $b=2$ exact PCA vs.\ $b=2$ randomized SVD), and (ii) stream (K vs.\ V). Head-to-head comparison against existing compressors is deferred to \S\ref{sec:benchmark-vs-turboquant}. All values are mean per-block reconstruction MSE on real captured tensors, averaged across all full-attention layers excluding layer~0 (RoPE-degenerate).

\paragraph{Configuration sweep (K stream).}
Table~\ref{tab:mse-k-sweep} reports the transition from the v1.2 baseline ($b=3$ exact PCA) through the bit-width reduction ($b=2$ exact) to the final v1.3 tier-1 configuration ($b=2$ randomized PCA with $r=D/2$).

\begin{table}[H]
\centering
\small
\caption{K-stream MSE across the three codec configurations at ctx=4096.}
\label{tab:mse-k-sweep}
\begin{tabular}{lrrrrr}
  \toprule
  Model & $b=3$ exact & $b=2$ exact & $b=2$ rsvd & $\frac{b=2e}{b=3e}$ & $\frac{b=2r}{b=3e}$ \\
  \midrule
  Qwen2.5-0.5B                  & $9.89\!\cdot\!10^{-1}$ & $1.10\!\cdot\!10^{0}$  & $1.11\!\cdot\!10^{0}$  & $1.11\times$ & $1.13\times$ \\
  Qwen3-0.6B                    & $1.35\!\cdot\!10^{0}$  & $1.44\!\cdot\!10^{0}$  & $1.46\!\cdot\!10^{0}$  & $1.07\times$ & $1.08\times$ \\
  Gemma-4-E2B                   & $3.78\!\cdot\!10^{-3}$ & $1.58\!\cdot\!10^{-3}$ & $1.58\!\cdot\!10^{-3}$ & $0.42\times$ & $0.42\times$ \\
  DeepSeek-R1-Distill           & $5.26\!\cdot\!10^{-1}$ & $5.83\!\cdot\!10^{-1}$ & $5.93\!\cdot\!10^{-1}$ & $1.11\times$ & $1.13\times$ \\
  GLM-Edge-1.5B                 & $7.65\!\cdot\!10^{-1}$ & $8.35\!\cdot\!10^{-1}$ & $8.49\!\cdot\!10^{-1}$ & $1.09\times$ & $1.11\times$ \\
  SmolLM2-1.7B                  & $6.73\!\cdot\!10^{-1}$ & $7.63\!\cdot\!10^{-1}$ & $7.80\!\cdot\!10^{-1}$ & $1.13\times$ & $1.16\times$ \\
  GLM-Edge-4B                   & $6.89\!\cdot\!10^{-1}$ & $7.56\!\cdot\!10^{-1}$ & $7.70\!\cdot\!10^{-1}$ & $1.10\times$ & $1.12\times$ \\
  \bottomrule
\end{tabular}
\end{table}

Two observations stand out. First, the transition from exact PCA to randomized PCA at $r=D/2$ adds only $1\%$--$2\%$ additional K-MSE on top of the $b=3 \to b=2$ bit-width change --- the randomized skeleton is essentially exact on the K stream because K is already low-rank (top-$10$ of $D=128$ captures $95\%$ variance on Qwen3 / DeepSeek). Second, Gemma-4 K-MSE is an outlier: $b=2$ \emph{improves} K quality relative to $b=3$ ($0.42\times$), because the $b=3$ Lloyd-Max codebook over-quantizes the near-zero residual on Gemma's high-head-dim K. This is a genuine RDO win, not a measurement artefact, and it is the only ACCEPT-without-tradeoff result in the matrix.

\paragraph{Configuration sweep (V stream).}
Table~\ref{tab:mse-v-sweep} gives the corresponding V-stream MSE trajectory. V is generally more tolerant than K under MSE distortion, except in SmolLM2's MHA head\_dim=64 regime (see \S\ref{sec:smollm2}).

\begin{table}[H]
\centering
\small
\caption{V-stream MSE across the three codec configurations at ctx=4096.}
\label{tab:mse-v-sweep}
\begin{tabular}{lrrrrr}
  \toprule
  Model & $b=3$ exact & $b=2$ exact & $b=2$ rsvd & $\frac{b=2e}{b=3e}$ & $\frac{b=2r}{b=3e}$ \\
  \midrule
  Qwen2.5-0.5B                  & $2.33\!\cdot\!10^{-1}$ & $2.64\!\cdot\!10^{-1}$ & $2.68\!\cdot\!10^{-1}$ & $1.13\times$ & $1.15\times$ \\
  Qwen3-0.6B                    & $4.34\!\cdot\!10^{0}$  & $4.73\!\cdot\!10^{0}$  & $5.29\!\cdot\!10^{0}$  & $1.09\times$ & $1.22\times$ \\
  Gemma-4-E2B                   & $2.81\!\cdot\!10^{-1}$ & $3.07\!\cdot\!10^{-1}$ & $3.03\!\cdot\!10^{-1}$ & $1.09\times$ & $1.08\times$ \\
  DeepSeek-R1-Distill           & $4.82\!\cdot\!10^{-1}$ & $5.26\!\cdot\!10^{-1}$ & $5.15\!\cdot\!10^{-1}$ & $1.09\times$ & $1.07\times$ \\
  GLM-Edge-1.5B                 & $3.64\!\cdot\!10^{-1}$ & $3.96\!\cdot\!10^{-1}$ & $4.10\!\cdot\!10^{-1}$ & $1.09\times$ & $1.13\times$ \\
  SmolLM2-1.7B                  & $2.68\!\cdot\!10^{-1}$ & $3.03\!\cdot\!10^{-1}$ & $\mathbf{4.33\!\cdot\!10^{-1}}$ & $1.13\times$ & $\mathbf{1.61\times}$ \\
  GLM-Edge-4B                   & $4.17\!\cdot\!10^{-1}$ & $4.55\!\cdot\!10^{-1}$ & $4.93\!\cdot\!10^{-1}$ & $1.09\times$ & $1.18\times$ \\
  \bottomrule
\end{tabular}
\end{table}

The SmolLM2 V-MSE inflation of $1.61\times$ is the REJECT signal that creates the tier-2 fallback documented in \S\ref{sec:smollm2}. On all other models V-MSE inflation stays $\leq 1.22\times$, within the MARGINAL band.

\paragraph{Inverse-RoPE K-MSE (tier-1.5).}
Table~\ref{tab:mse-rope} reports the full per-model pre/post-RoPE K-MSE numbers underlying the tier-1.5 decision. Layer~0 is excluded (RoPE is identity at position $0$).

\begin{table}[H]
\centering
\small
\caption{K-stream MSE under the inverse-RoPE preprocessor versus post-RoPE (ctx=4096). $\downarrow$ means MSE decreased; $\uparrow$ means it increased.}
\label{tab:mse-rope}
\begin{tabular}{lrrrl}
  \toprule
  Model & post-RoPE K-MSE & pre-RoPE K-MSE & ratio & Verdict \\
  \midrule
  Qwen2.5-0.5B                  & $1.33\!\cdot\!10^{0}$  & $6.50\!\cdot\!10^{-1}$ & $0.49\times\downarrow$  & ACCEPT \\
  Qwen3-0.6B                    & $2.36\!\cdot\!10^{0}$  & $2.03\!\cdot\!10^{0}$  & $0.86\times\downarrow$  & ACCEPT \\
  DeepSeek-R1-Distill           & $6.53\!\cdot\!10^{-1}$ & $3.76\!\cdot\!10^{-1}$ & $0.58\times\downarrow$  & ACCEPT \\
  SmolLM2-1.7B                  & $1.63\!\cdot\!10^{0}$  & $1.50\!\cdot\!10^{0}$  & $0.92\times\downarrow$  & MARGINAL \\
  Gemma-4-E2B                   & $1.58\!\cdot\!10^{-3}$ & $1.51\!\cdot\!10^{-3}$ & $0.95\times\downarrow$  & REJECT (bytes worsen) \\
  GLM-Edge-1.5B                 & $9.52\!\cdot\!10^{-1}$ & $1.07\!\cdot\!10^{0}$  & $1.13\times\uparrow$    & REJECT \\
  GLM-Edge-4B                   & $9.00\!\cdot\!10^{-1}$ & $1.01\!\cdot\!10^{0}$  & $1.12\times\uparrow$    & REJECT \\
  \bottomrule
\end{tabular}
\end{table}

The clean bifurcation (halfsplit-RoPE-without-QK-norm models accept; all others reject) is the foundation of the tier-1.5 capability table in \S\ref{sec:benchmark}.

\paragraph{Per-layer K-MSE distribution.}
Beyond the aggregate mean, the per-layer MSE distribution is informative because attention-critical layers may sit at distribution tails. Table~\ref{tab:mse-perlayer} gives the Qwen3-0.6B K-stream quantiles under tier-1.

\begin{table}[H]
\centering
\small
\caption{Per-layer K-MSE distribution on Qwen3-0.6B under tier-1, over 27 full-attention layers (L0 excluded).}
\label{tab:mse-perlayer}
\begin{tabular}{lr}
  \toprule
  Statistic & K-MSE \\
  \midrule
  min (layer 13) & $8.08\!\cdot\!10^{-1}$ \\
  25th percentile & $9.70\!\cdot\!10^{-1}$ \\
  median          & $1.14\!\cdot\!10^{0}$  \\
  75th percentile & $1.89\!\cdot\!10^{0}$  \\
  max (layer 4)   & $3.15\!\cdot\!10^{0}$  \\
  mean            & $1.46\!\cdot\!10^{0}$  \\
  \bottomrule
\end{tabular}
\end{table}

The max-to-min spread across layers is $3.9\times$. Because each layer fits its own low-rank skeleton (the RSVD is per block, not global), worst-case layers do not diverge further than this spread --- a property we will contrast with data-oblivious codecs in \S\ref{sec:benchmark-vs-turboquant}.

\subsection{Pareto summary across configurations}
\label{sec:pareto}

Ranking \textsc{KakeyaTurbo} configurations on the $(\text{ratio}, \text{quality})$ plane:
\begin{itemize}[leftmargin=*]
  \item Exact PCA with $b=3$: lowest ratio ($\sim 3\times$), highest quality (ACCEPT by construction).
  \item Tier-1 ($b=2$ exact): mid ratio ($\sim 4.5\times$), MARGINAL quality.
  \item \textsc{Tier-1} ($b=2$, $r=D/2$ randomized SVD): high ratio ($5.4\times$--$6.7\times$), MARGINAL quality on six of seven models, ACCEPT on Gemma-4.
  \item \textsc{Tier-1.5} (tier-1 + inverse-RoPE): moves the Qwen / DeepSeek family further by $+8\%$ to $+9\%$ ratio at no quality cost.
  \item SmolLM2 is the only model where the full tier-1 ratio cannot be reached within the ACCEPT MSE band (see \S\ref{sec:smollm2}).
\end{itemize}

\subsection{Flagship projections at 2026-Q2 scale}

None of the April-2026 flagship open-source models (Qwen3-235B-A22B, DeepSeek-V3.1, Kimi-K2-Instruct, GLM-4.6, MiniMax-M2) is loadable on our 15\,GB-RAM VM; their bf16 weights alone are $458$--$2000$\,GB. We therefore report byte-exact analytical extrapolation from architecturally-identical small-sibling proxies that \emph{are} measured (Table~\ref{tab:flagship}).

\begin{table}[H]
\centering
\small
\caption{Byte-exact extrapolation to 2026-Q2 flagship open-source models. Per-vector bytes come from the measured small-sibling proxy; $L$ and $n_{kv} \cdot \text{hd}$ come from the flagship config.}
\label{tab:flagship}
\begin{tabular}{lrrrrr}
  \toprule
  Flagship & Params & $L$ & hd & tier-1 & tier-1.5 \\
  \midrule
  Qwen3-235B-A22B           & 235\,B & 94 & 128 & $6.53\times$ & $\mathbf{7.13\times}$ \\
  DeepSeek-V3.1 (decomp)    & 671\,B & 61 & 192 & $5.92\times$ & $6.76\times$ \\
  Kimi-K2-Instruct (decomp) & $1{,}000$\,B & 61 & 192 & $5.92\times$ & $6.76\times$ \\
  GLM-4.6                   & 355\,B & 92 & 128 & $5.84\times$ & --- \\
  MiniMax-M2                & 229\,B & 62 & 128 & $5.92\times$ & $6.76\times$ \\
  \bottomrule
\end{tabular}
\end{table}

DeepSeek-V3.1 and Kimi-K2 use Multi-Latent Attention (MLA); the reported ratios are on the decompressed K/V (what the attention computation sees). The additional gain from applying \textsc{KakeyaTurbo} to the stored MLA latent is an open question.

\subsection{Head-to-head comparison with TurboQuant turbo3}
\label{sec:benchmark-vs-turboquant}

This subsection consolidates all direct comparisons against the strongest existing post-hoc KV compression baseline, TurboQuant turbo3~\cite{turboquant}, on the exact same captured tensors. We keep the comparison in one place so the rest of the paper can focus on \textsc{KakeyaTurbo}'s own algorithm and RD boundary.

\paragraph{Compression ratio (tier-1, ctx=4\,096 and ctx=128\,k).}
Table~\ref{tab:bench-ratio-vs-turbo3} gives the ratio comparison on the seven benchmark models.

\begin{table}[H]
\centering
\small
\caption{Compression ratio vs.\ bf16 baseline: \textsc{KakeyaTurbo} tier-1 ($b=2$, $r=D/2$) against TurboQuant turbo3 ($3$ bits/coord on K+V + fp32 K-norm), measured at ctx=4\,096 and byte-exactly extrapolated to ctx=128\,k.}
\label{tab:bench-ratio-vs-turbo3}
\begin{tabular}{lrrrrr}
  \toprule
  Model & \multicolumn{2}{c}{ctx = 4\,096} & \multicolumn{2}{c}{ctx = 128k} & $\Delta$ \\
  \cmidrule(lr){2-3} \cmidrule(lr){4-5}
  & \textsc{Turbo} & turbo3 & \textsc{Turbo} & turbo3 & \\
  \midrule
  Qwen2.5-0.5B          & $5.40\times$ & $4.92\times$ & $5.46\times$ & $4.92\times$ & $+11\%$ \\
  Qwen3-0.6B            & $\mathbf{6.61\times}$ & $5.12\times$ & $\mathbf{6.65\times}$ & $5.12\times$ & $\mathbf{+30\%}$ \\
  Gemma-4-E2B (FA only) & $\mathbf{6.32\times}$ & $5.28\times$ & $\mathbf{6.72\times}$ & $5.28\times$ & $\mathbf{+27\%}$ \\
  DeepSeek-R1-Distill   & $5.98\times$ & $5.12\times$ & $6.11\times$ & $5.12\times$ & $+19\%$ \\
  GLM-Edge-1.5B         & $5.85\times$ & $5.12\times$ & $5.91\times$ & $5.12\times$ & $+15\%$ \\
  SmolLM2-1.7B          & $5.37\times$ & $4.92\times$ & $5.37\times$ & $4.92\times$ & $+9\%$ \\
  GLM-Edge-4B           & $5.83\times$ & $5.12\times$ & $5.88\times$ & $5.12\times$ & $+15\%$ \\
  \midrule
  \textbf{Average}      & \textbf{5.91}$\times$ & 5.08$\times$ & \textbf{6.01}$\times$ & 5.08$\times$ & \textbf{+18\%} \\
  \bottomrule
\end{tabular}
\end{table}

Six of seven models achieve a net positive ratio margin over turbo3 at ACCEPT-level quality; SmolLM2 reaches tier-1 ratio only at MARGINAL V-MSE (\S\ref{sec:smollm2}). Turning on the tier-1.5 inverse-RoPE preprocessor takes the Qwen / DeepSeek family further to $+30\%$ to $+39\%$ over turbo3 (cross-referencing Table~\ref{tab:rope}).

\paragraph{Reconstruction MSE.}
Per-layer mean K-stream MSE on the identical captured tensors, summarised in Table~\ref{tab:mse-vs-turbo3}, reveals that the ratio advantage is not traded for quality --- on the Qwen / DeepSeek family the gap is large enough that \textsc{KakeyaTurbo} is strictly dominant.

\begin{table}[H]
\centering
\small
\caption{Head-to-head K-stream MSE on identical captured tensors (ctx=4\,096). Ratio $>1$ means \textsc{KakeyaTurbo} has lower MSE. turbo3 MSE numbers come from our prior comparison~\cite{kakeya-kv}.}
\label{tab:mse-vs-turbo3}
\begin{tabular}{lrrrl}
  \toprule
  Model & \textsc{Turbo} K-MSE & turbo3 K-MSE & turbo3 / \textsc{Turbo} & Interpretation \\
  \midrule
  Qwen2.5-0.5B                  & $1.11\!\cdot\!10^{0}$  & $2.07\!\cdot\!10^{2}$  & $\mathbf{186\times}$   & \textsc{KakeyaTurbo} much lower \\
  Qwen3-0.6B                    & $1.46\!\cdot\!10^{0}$  & $8.99\!\cdot\!10^{1}$  & $\mathbf{62\times}$    & \textsc{KakeyaTurbo} much lower \\
  DeepSeek-R1-Distill           & $5.93\!\cdot\!10^{-1}$ & $1.44\!\cdot\!10^{3}$  & $\mathbf{2428\times}$  & \textsc{KakeyaTurbo} much lower \\
  Gemma-4-E2B                   & $1.58\!\cdot\!10^{-3}$ & $2.96\!\cdot\!10^{-3}$ & $1.87\times$           & \textsc{KakeyaTurbo} lower \\
  GLM-Edge-1.5B                 & $8.49\!\cdot\!10^{-1}$ & $2.81\!\cdot\!10^{-1}$ & $0.33\times$           & turbo3 lower \\
  GLM-Edge-4B                   & $7.70\!\cdot\!10^{-1}$ & $1.48\!\cdot\!10^{-1}$ & $0.19\times$           & turbo3 lower \\
  \bottomrule
\end{tabular}
\end{table}

On the Qwen / DeepSeek family (halfsplit RoPE without QK-norm) \textsc{KakeyaTurbo}'s K-MSE is $62\times$ to $2428\times$ lower than turbo3, consistent with the asymmetric-rescue configurations the TurboQuant documentation recommends for these architectures. On GLM-Edge (adjacent-pairs RoPE with QK-norm) turbo3 has lower K-MSE because QK-norm already norm-stabilises the K representation and the Gaussianised scalar quantiser handles it cleanly. Gemma-4 sits in between, with \textsc{KakeyaTurbo} K-MSE $1.87\times$ lower due to Gemma's high-head-dim low-rank K structure.

\paragraph{Cross-layer variance.}
A secondary but informative comparison is the spread of K-MSE across layers. \textsc{KakeyaTurbo}'s per-block PCA re-fits each layer's low-rank structure, so the worst-case layer does not diverge from the best by more than $3.9\times$ on Qwen3-0.6B (Table~\ref{tab:mse-perlayer}). turbo3's fixed rotation produces a $\sim 200\times$ cross-layer spread on the same tensor set, reflecting the fact that a data-oblivious rotation cannot adapt to per-layer anisotropy.

\paragraph{Summary.}
At comparable byte budget and same captured-tensor quality metric, \textsc{KakeyaTurbo} tier-1 strictly dominates turbo3 on six of seven benchmark models; tier-1.5 dominates on Qwen / DeepSeek by a further $+30\%$ to $+39\%$ in ratio; on GLM-Edge, turbo3 has an MSE advantage that reflects its natural fit with QK-norm architectures. Flagship extrapolation (Table~\ref{tab:flagship}) projects this ordering to $128$k context on Qwen3-235B, DeepSeek-V3.1, Kimi-K2, GLM-4.6 and MiniMax-M2.

\section{Related work}
\label{sec:related}

\paragraph{KV cache compression.} H\textsubscript{2}O~\cite{h2o} and Scissorhands~\cite{scissorhands} propose token eviction based on attention-score heavy hitters. StreamingLLM~\cite{streaming-llm} observes the attention-sink phenomenon and preserves the first few tokens exact. KIVI~\cite{kivi} is a per-channel quantization scheme for K (asymmetric quantization per channel) and per-token for V. TurboQuant~\cite{turboquant} is a data-oblivious WHT + Lloyd--Max scalar quantizer. We compare against turbo3 head-to-head on the same captured tensors in \S\ref{sec:benchmark}.

\paragraph{Low-rank KV methods.} Kakeya-KV~\cite{kakeya-kv} is our prior work and the direct ancestor of the skeleton construction here. DMC~\cite{dmc} and LESS~\cite{less} learn low-rank projections during training. Our approach is post-hoc and training-free.

\paragraph{Randomized numerical linear algebra.} The Halko--Martinsson--Tropp randomized SVD~\cite{halko-martinsson-tropp} is a textbook algorithm; we use it as a drop-in for exact PCA and observe that the implicit rank cap is the structural win. Frequent Directions~\cite{frequent-directions} is a related streaming PCA sketch we defer to future work.

\paragraph{Kakeya geometry and harmonic analysis.} The Kakeya problem originates in~\cite{kakeya-problem}. Its deep role in harmonic analysis began with Fefferman~\cite{fefferman-ball} (ball multiplier), continued through Bourgain's arithmetic-combinatorial bounds~\cite{bourgain-kakeya}, Wolff's geometric hairbrush argument~\cite{wolff-kakeya}, and Katz--{\L}aba--Tao's $\R^3$ improvements~\cite{katz-laba-tao}. The finite-field analogue was settled by Dvir's polynomial method~\cite{dvir-finite-field}; Bennett--Carbery--Tao's multilinear Kakeya conjecture~\cite{bennett-carbery-tao} and Guth's endpoint resolution~\cite{guth-endpoint} introduced the tools eventually used by Wang and Zahl~\cite{wang-zahl} to resolve the Euclidean three-dimensional case. \emph{Our use of this chain is formal, not intuitional}, as documented in \S\ref{sec:kbl-chain}: multilinear Kakeya is equivalent to Brascamp--Lieb~\cite{bennett-carbery-christ-tao,carbery-valdimarsson}, which underpins Tropp's matrix concentration inequalities~\cite{tropp-matrix-chernoff}, which underpin the Halko--Martinsson--Tropp RSVD error bound~\cite{halko-martinsson-tropp}. The codec skeleton is therefore a literature-traceable shallow instance of the same toolchain used at its deep end in~\cite{wang-zahl}. Proposition~\ref{prop:rsvd-skeleton} makes the quantitative correspondence precise; the structural parallels (power iteration as multiscale induction, Gaussian probing as direction enumeration, singular-value distribution as Frostman measure) and the three rigorous disclaimers are in \S\ref{sec:kbl-chain}.

\section{Limitations and future work}
\label{sec:limitations}

\paragraph{Quality measurement.} We report reconstruction MSE, not perplexity or downstream-task accuracy. A full PPL sweep across LongBench/RULER is deferred to future work; early checks on Gemma-4-E2B at 2k context show greedy-decode token-level agreement with the bf16 baseline for the first 12 tokens. End-to-end quality at 128k context on flagship models remains unverified.

\paragraph{Flagship measurements.} All flagship numbers are analytical extrapolation from measured proxies. Validation on $\geq 500$\,GB-RAM hardware is pending.

\paragraph{MLA models.} DeepSeek-V3.1 and Kimi-K2 use MLA; our reported numbers compress the decompressed K/V, not the stored latent. Applying the codec to the MLA latent is a v1.4 open item.

\paragraph{GLM-family inverse-RoPE.} GLM-Edge uses adjacent-pairs RoPE + QK-norm; our halfsplit inverse is incompatible. A GLM-correct inverse rotation is tracked as v1.3.1.

\paragraph{Attention-aware extensions.} The \texttt{weights: \&[f32]} parameter in our API is a reserved socket for attention-score-driven importance, but is not yet wired. The v1.4 roadmap (public in the repository) proposes closing this loop by piping H\textsubscript{2}O-style prefill attention scores into the weighted PCA fit.

\paragraph{SmolLM2-class models.} MHA with head dim $64$ and flat V spectra defeats rank truncation; the Pareto gap observed on SmolLM2 will likely appear on any future model with similar structure. We ship SmolLM2 as a tier-2 capability that requires an explicit quality-versus-bytes trade-off at the deployment policy layer.

\section{Conclusion}
\label{sec:conclusion}

We presented \textsc{KakeyaTurbo}, a training-free, three-stage KV cache compression codec, and established its rate--distortion boundary.

On the \emph{algorithm} side, the three stages --- weighted randomized SVD with rank cap $r = D/2$, per-architecture inverse-RoPE preprocessing on the key stream, and 2-bit Lloyd--Max residual quantization after WHT --- are the separable optimisers of three nested subproblems of a single rate--distortion objective, not a pipeline of independent heuristics. The codec is parameterised entirely by a compile-time distortion metric and a per-vector weight, and runs in the full LLM inference pipeline (prefill + token-by-token decode + paged attention) with a block-streaming mode whose amortised overhead is under $10\,\mu$s per layer per decode step on H100.

On the \emph{RD boundary} side, the randomized-SVD skeleton sits at the shallow end of an explicit four-step dependency chain whose deep end is used by Wang and Zahl in their recent resolution of the three-dimensional Kakeya conjecture. The chain yields an unconditional upper bound (Proposition~\ref{prop:rsvd-skeleton}) and a conjectural lower bound (Proposition~\ref{prop:kakeya-rd}) that closes unconditionally at $D = 3$ via Wang--Zahl. The rank cap $r = D/2$ is thereby derived as a Pareto-optimal point on the RD curve.

The supporting benchmark shows that \textsc{KakeyaTurbo} achieves $5.4\times$--$6.7\times$ compression against the bf16 baseline on six of seven open-source models at ACCEPT quality, with analytical projections to $7.1\times$ on flagship Qwen3-235B. The head-to-head comparison against the strongest prior compressor (\S\ref{sec:benchmark-vs-turboquant}) confirms the ordering. Open items are limited and specific: end-to-end perplexity validation, GLM-Edge inverse-RoPE (adjacent-pairs + QK-norm), MLA-latent compression, and the attention-aware extension; each has a concrete roadmap in the public repository.

\paragraph{Reproducibility.} All code, unit tests, integration tests, property-based tests, real-measurement JSON, and ablation decision documents are released under Apache-2.0. Repository:
\begin{center}
\url{https://github.com/FluffyAIcode/LLM-KV--Cache-compress}
\end{center}
The v1.3 release state lives on branch \texttt{cursor/v1-3-rsvd-rope-aware-12f5} (Pull Request \#11); the v1.4 roadmap is in \texttt{reports/v1\_3\_rsvd\_rope/V1\_4\_V1\_5\_ROADMAP.md}.

\begin{thebibliography}{99}

\bibitem{h2o}
Zhang, Z., Sheng, Y., Zhou, T., Chen, T., Zheng, L., Cai, R., Song, Z., Tian, Y., R\'e, C., Barrett, C., Wang, Z., Chen, B.
\newblock {H$_2$O: Heavy-hitter oracle for efficient generative inference of large language models.}
\newblock \emph{NeurIPS}, 2023.

\bibitem{scissorhands}
Liu, Z., Desai, A., Liao, F., Wang, W., Xie, V., Xu, Z., Kyrillidis, A., Shrivastava, A.
\newblock {Scissorhands: Exploiting the persistence of importance hypothesis for LLM KV cache compression at test time.}
\newblock \emph{NeurIPS}, 2023.

\bibitem{deepseekv3}
DeepSeek-AI.
\newblock {DeepSeek-V3 Technical Report.}
\newblock \emph{arXiv preprint arXiv:2412.19437}, 2024.

\bibitem{turboquant}
TheTom.
\newblock {TurboQuant+: Gaussian-bounded scalar quantization for KV cache.}
\newblock \emph{GitHub repository, \url{https://github.com/TheTom/turboquant_plus}}, 2024.

\bibitem{kakeya-kv}
KakeyaKV Contributors.
\newblock {Kakeya KV cache compression.}
\newblock \emph{\url{https://github.com/FluffyAIcode/LLM-KV--Cache-compress}}, 2026. (v1.0 and v1.2, this repository.)

\bibitem{kakeya-problem}
Kakeya, S.
\newblock {Some problems on maxima and minima regarding ovals.}
\newblock \emph{Tohoku Science Reports}, 6:71--88, 1917.

\bibitem{fefferman-ball}
Fefferman, C.
\newblock {The multiplier problem for the ball.}
\newblock \emph{Annals of Mathematics}, 94(2):330--336, 1971.

\bibitem{bourgain-kakeya}
Bourgain, J.
\newblock {Besicovitch type maximal operators and applications to Fourier analysis.}
\newblock \emph{Geometric and Functional Analysis}, 1(2):147--187, 1991.

\bibitem{wolff-kakeya}
Wolff, T.
\newblock {An improved bound for Kakeya type maximal functions.}
\newblock \emph{Revista Matem\'atica Iberoamericana}, 11(3):651--674, 1995.

\bibitem{katz-laba-tao}
Katz, N.~H., {\L}aba, I., Tao, T.
\newblock {An improved bound on the Minkowski dimension of Besicovitch sets in $\R^3$.}
\newblock \emph{Annals of Mathematics}, 152(2):383--446, 2000.

\bibitem{dvir-finite-field}
Dvir, Z.
\newblock {On the size of Kakeya sets in finite fields.}
\newblock \emph{Journal of the American Mathematical Society}, 22(4):1093--1097, 2009.

\bibitem{bennett-carbery-tao}
Bennett, J., Carbery, A., Tao, T.
\newblock {On the multilinear restriction and Kakeya conjectures.}
\newblock \emph{Acta Mathematica}, 196(2):261--302, 2006.

\bibitem{guth-endpoint}
Guth, L.
\newblock {The endpoint case of the Bennett-Carbery-Tao multilinear Kakeya conjecture.}
\newblock \emph{Acta Mathematica}, 205(2):263--286, 2010.

\bibitem{bennett-carbery-christ-tao}
Bennett, J., Carbery, A., Christ, M., Tao, T.
\newblock {The Brascamp--Lieb inequalities: finiteness, structure and extremals.}
\newblock \emph{Geometric and Functional Analysis}, 17(5):1343--1415, 2008.

\bibitem{carbery-valdimarsson}
Carbery, A., Valdimarsson, S.~I.
\newblock {The endpoint multilinear Kakeya theorem via the Borsuk--Ulam theorem.}
\newblock \emph{Journal of Functional Analysis}, 264(7):1643--1663, 2013.

\bibitem{tropp-matrix-chernoff}
Tropp, J.~A.
\newblock {User-friendly tail bounds for sums of random matrices.}
\newblock \emph{Foundations of Computational Mathematics}, 12(4):389--434, 2012.

\bibitem{wang-zahl}
Wang, H., Zahl, J.
\newblock {Volume estimates for unions of convex sets, and the Kakeya set conjecture in three dimensions.}
\newblock \emph{arXiv preprint arXiv:2502.17655}, 2025.

\bibitem{halko-martinsson-tropp}
Halko, N., Martinsson, P.-G., Tropp, J.~A.
\newblock {Finding structure with randomness: Probabilistic algorithms for constructing approximate matrix decompositions.}
\newblock \emph{SIAM Review}, 53(2):217--288, 2011.

\bibitem{rope}
Su, J., Lu, Y., Pan, S., Murtadha, A., Wen, B., Liu, Y.
\newblock {RoFormer: Enhanced transformer with rotary position embedding.}
\newblock \emph{Neurocomputing}, 568:127063, 2024.

\bibitem{streaming-llm}
Xiao, G., Tian, Y., Chen, B., Han, S., Lewis, M.
\newblock {Efficient streaming language models with attention sinks.}
\newblock \emph{ICLR}, 2024.

\bibitem{kivi}
Liu, Z., Yuan, J., Jin, H., Zhong, S., Xu, Z., Braverman, V., Chen, B., Hu, X.
\newblock {KIVI: A tuning-free asymmetric 2bit quantization for KV cache.}
\newblock \emph{ICML}, 2024.

\bibitem{dmc}
Nawrot, P., Lancucki, A., Chochowski, M., Tarjan, D., Ponti, E.~M.
\newblock {Dynamic memory compression: Retrofitting LLMs for accelerated inference.}
\newblock \emph{arXiv preprint arXiv:2403.09636}, 2024.

\bibitem{less}
Dong, H., Yang, X., Zhang, Z., Wang, Z., Chi, Y., Chen, B.
\newblock {Get more with LESS: Synthesizing recurrence with KV cache compression for efficient LLM inference.}
\newblock \emph{ICML}, 2024.

\bibitem{frequent-directions}
Liberty, E.
\newblock {Simple and deterministic matrix sketching.}
\newblock \emph{KDD}, 2013.

\end{thebibliography}

\appendix

\section{Complete benchmark protocol}
\label{app:protocol}

All benchmarks in this paper are reproducible via the following one-line invocation after cloning the repository and downloading a model:
\begin{verbatim}
python3 benchmarks/kakeyaturbo_v1_2_real_bench.py \
    --model-path models/<model-dir> \
    --model-name <model-name> \
    --context-tokens 4096 \
    --bit-width 2 \
    --pca-method randomized \
    --out-dir reports/v1_3_rsvd_rope/bench/<model>_rsvd_half/ctx_4096
\end{verbatim}
The inverse-RoPE preprocessor adds:
\begin{verbatim}
python3 benchmarks/rope_aware_k_poc.py \
    --model-dir <model-dir> --model-name <model-name> \
    --ctx 4096 --rope-pairing halfsplit \
    --out-dir reports/v1_3_rsvd_rope/rope_poc/<model-name>
\end{verbatim}

Each run produces a JSON manifest per layer plus an aggregate \verb|summary.json|. All manifests are committed to the repository under \verb|reports/v1_3_rsvd_rope/|.

\section{Codec parameter defaults}
\label{app:params}

The reference v1.3 tier-1 configuration is:
\begin{center}
\begin{tabular}{ll}
  \toprule
  Parameter & Default value \\
  \midrule
  \texttt{bit\_width} & 2 \\
  \texttt{variance\_ratio} & 0.95 \\
  \texttt{k} (K-means centres) & 16 \\
  \texttt{block\_size} & 512 \\
  \texttt{kmeans\_max\_iter} & 32 \\
  \texttt{share\_basis} (V stream) & true \\
  \texttt{share\_basis} (K stream) & false \\
  \bottomrule
\end{tabular}
\end{center}

The randomized-SVD parameters are:
\begin{center}
\begin{tabular}{ll}
  \toprule
  Parameter & Default value \\
  \midrule
  \texttt{target\_rank} & $D/2$ \\
  \texttt{oversample} ($p$) & 8 \\
  \texttt{power\_iters} ($q$) & 2 \\
  \texttt{seed\_offset} & \texttt{0x9E37\_79B9\_7F4A\_7C15} \\
  \bottomrule
\end{tabular}
\end{center}

Tier-1.5 additionally applies the inverse-RoPE preprocessor to the K stream before the codec sees it, using the model's \texttt{rope\_theta} and halfsplit pairing convention.

\section{Qualitative ablation trail}
\label{app:ablation}

The v1.3 configuration is the endpoint of seven ablations run over the project's lifetime. For completeness, the full decision chain is:
\begin{enumerate}[leftmargin=*]
  \item \textbf{v1.0 $\to$ v1.2 A}: bf16 storage of skeleton (halves skeleton bytes, no quality loss). ACCEPT, shipped.
  \item \textbf{v1.2 B'}: V-stream layer-pooled PCA basis. ACCEPT (V only), shipped. Rejected on K due to family-dependent divergence.
  \item \textbf{K-side $d_{\text{eff}}$ truncation} + outlier compensation: every $0.05$ drop in \texttt{variance\_ratio} doubles K MSE; outliers recover $\sim 25\%$. \textbf{REJECT}.
  \item \textbf{Block size doubling} (512 $\to$ 1024): MSE inflation $1.14\times$ on average with $16\%$ byte savings. \textbf{MARGINAL}, not shipped as default.
  \item \textbf{Cross-layer basis sharing} on K: $2\times$--$15\times$ MSE inflation on Qwen/DeepSeek family, ACCEPT on Gemma/GLM-Edge. \textbf{REJECT as default, accept as per-family knob.}
  \item \textbf{Bit width sweep} ($b=3 \to b=2$): $+4\%$ to $+23\%$ ratio, MSE inflation $1.07\times$--$1.13\times$. \textbf{ACCEPT, new default.}
  \item \textbf{Randomized SVD with rank cap $r = D/2$}: combined with $b=2$, pushes tier-1 compression ratio to $5.4\times$--$6.7\times$ on 6/7 models at ACCEPT quality. \textbf{ACCEPT, new default.}
  \item \textbf{Inverse-RoPE K preprocessor}: $14\%$--$51\%$ K MSE drop on halfsplit-RoPE family, no-op or harmful elsewhere. \textbf{ACCEPT as per-family knob.}
\end{enumerate}
The full narrative is in the repository's \verb|reports/v1_3_rsvd_rope/DECISION.md|.

\end{document}
